\section{Open Questions}
\label{sec:open-questions}

These five questions are genuinely open after the replication effort. Each is grounded in specific gaps in the paper (or in this reproduction's scope) and carries concrete next-steps executable on the current free-endpoint compute (uicgpu Geant4-DNA, CPU LEM-IV/MKM, PIDE database).

\subsection*{Q1. Independent MC cross-code validation}
\textbf{Question.} How well do the Kundr\'at/PARTRAC analytical formulas agree with an independent MCTS code (Geant4-DNA or TRAX-CHEM) run on the same PARTRAC target geometry (10~$\mu$m spherical lymphocyte nucleus, 6.6~Gbp, G0/G1) for the same ion/energy grid, across the full published LET range and specifically at the SSB dip (5--20~keV/$\mu$m) and the DSB peak (100--500~keV/$\mu$m)?

\textbf{Basis.} The paper validates the formulas only against its own PARTRAC MCTS outputs; refs~28--31 note ongoing Geant4-DNA implementations of comparable physics but no cross-code residuals are published. Refs~12--14 (Friedland-group PARTRAC benchmarking) show that PARTRAC-vs-Geant4-DNA yield differences at matched geometry can reach a factor of 2 in DSB cluster fractions. No such cross-code validation of \emph{these} formulas exists.

\textbf{Next steps.} Run Geant4-DNA (available on \texttt{uicgpu}) with the \texttt{dna} physics list on a spherical water target with implanted geometric-DNA scoring at 6.6~Gbp equivalent for protons (0.5--100~MeV/u, 6 energies) and $^{4}$He (1--100~MeV/u, 6 energies); score SB/SSB/DSB per Gy per Gbp; compute residual $(\text{formula}-\text{MC})/\text{formula}$ at each LET; report $\chi^2$ per damage class.

\subsection*{Q2. LET-range extrapolation safety}
\textbf{Question.} Are the formulas safe to extrapolate outside their 0.25--512~MeV/u training range, in particular for ultra-low LET ($\ll 0.3$~keV/$\mu$m) and ultra-high LET ($>1000$~keV/$\mu$m for $^{16}$O/$^{20}$Ne near/past the Bragg peak)?

\textbf{Basis.} Eq.~(2) contains a logistic overkill turnover but no built-in stopping-power ceiling. The paper explicitly declares the fit range and says nothing about extrapolation safety. For SB (Eq.~1), the low-LET plateau $p_1$ is hand-fixed across ions; whether this survives below the fit floor is untested.

\textbf{Next steps.} (a) Evaluate at LET~$=0.01, 0.05, 0.1$~keV/$\mu$m and check whether SB/SSB drop unphysically below the electron/gamma limit ($\sim$130 SB/Gy/Gbp). (b) Evaluate at LET~$=2000, 5000$~keV/$\mu$m and inspect large-LET asymptotics of Eq.~(2), which for typical $p_3\!\sim\!1.5$, $p_5\!\sim\!1.2$ goes as $\mathrm{LET}^{p_3-p_5}\!\sim\!\mathrm{LET}^{0.3}$ (keeps growing, likely unphysical). (c) Publish a documented ``safe range'' addendum.

\subsection*{Q3. Chromatin-heterogeneity extension}
\textbf{Question.} Can the analytical fits be re-parameterized to carry a chromatin-configuration variable (condensation fraction, heterochromatin/euchromatin ratio, S-phase replication-fork density) so they apply beyond the PARTRAC G0/G1 spherical-lymphocyte target?

\textbf{Basis.} The paper's target model is a single G0/G1 lymphocyte-nucleus configuration; all 540 fit parameters absorb the chromatin state implicitly. Damage yields for mitotic chromatin ($\sim 5\times$ more compact) or S-phase differ experimentally by up to 2--3$\times$. The formulas as-published cannot be dialed to cell state.

\textbf{Next steps.} Run PARTRAC (or Geant4-DNA-Chem, which we can) with three target geometries (G0/G1 spherical, condensed mitotic, S-phase with replication-fork nodes) at three ion/energy points each. Fit an added exponential-envelope factor $\exp(-\alpha \cdot \rho_{\text{chromatin}})$ to $(p_1, p_2, p_4)$ in Eq.~(1) and to numerator/denominator scales in Eq.~(2). If a single 3-parameter chromatin correction absorbs the geometry variation to $<20\%$, publish it as an extension formula.

\subsection*{Q4. Mapping to cellular survival (LEM/MKM/RMF)}
\textbf{Question.} How do the DSB / DSB-cluster / DSB-site yields from Eq.~(2) translate into cellular survival predictions when fed into an established biophysical survival model (LEM-IV, MKM, RMF, or PARTRAC's own repair chain), and does the analytical shortcut give the same survival curves as a full MCTS+repair chain?

\textbf{Basis.} The paper stops at DSB-yield output; downstream survival is out of scope. But the raison d'\^etre of these formulas is treatment planning, where survival --- not DSB yield --- is the endpoint. LEM-IV in particular is sensitive to DSB cluster fraction, so the shape of Fig.~4 matters more than Fig.~3 for RBE prediction.

\textbf{Next steps.} Wire the analytical yields into LEM-IV (public C++ implementation) and MKM (public Python implementations) as the DSB-input layer; compute survival curves $S(D)$ for 2~Gy fractions across the 9 ions at 4 LET points each; compare against experimental survival data compiled in the PIDE database; report RBE${}_{10\%}$ residuals. All free-endpoint (CPU).

\subsection*{Q5. Robustness to alternative direct/indirect chemistry}
\textbf{Question.} Are the direct-vs-indirect damage split assumptions (linear SB probability 0$\rightarrow$1 from 5--37.5~eV; 65\% breakage per $\bullet$OH attack) robust, and how sensitive are the analytical formulas to reasonable alternative assumptions used by other track-structure groups (KURBUC's stochastic OH scavenging; Geant4-DNA-Chem's energy-dependent OH yields)?

\textbf{Basis.} The paper commits to fixed direct/indirect probabilities inherited from earlier PARTRAC papers (Friedland refs~14--16). If the direct/indirect split changes, the shape of the Lorentzian dip in Eq.~(1) (driven by the indirect channel at intermediate LET) will change. This sensitivity is not analyzed anywhere in the paper.

\textbf{Next steps.} Re-fit Eq.~(1) parameters ($p_2, p_3, p_4, p_5$ for the SSB indirect channel) under two alternative indirect-chemistry assumptions: (a) OH breakage probability scaled by $0.75\times$ and $1.25\times$ uniformly; (b) energy-dependent scaling matching Geant4-DNA-Chem 11.2 yields. Report parameter shift magnitude and yield residual at LET~$=1, 10, 50, 200$~keV/$\mu$m. If shifts are $<20\%$ for the plateau and $<2\times$ for the dip depth, the formulas are robust; else publish the sensitivity band as a fit systematic.
